% Part: many-valued-logic
% Chapter: infinite-valued-logics
% Section: introduction

\documentclass[../../../include/open-logic-section]{subfiles}

\begin{document}

\olfileid{mvl}{inf}{int}

\olsection{مقدمه}

تعداد ارزش‌های صدق یک ماتریس لازم نیست متناهی باشد. یک انتخاب آشکار
برای مجموعه‌ای نامتناهی از ارزش‌های صدق، مجموعهٔ اعداد گویای میان~$0$
و~$1$، یعنی~$V_\infty = [0,1] \cap \Rat$ است؛ به بیان دقیق‌تر،
\begin{align*}
    V_\infty & = \Setabs{\frac{n}{m}}{n,m \in \Nat \text{ و } m>0
      \text{ و } n\le m}.
\intertext{هنگام بررسی این مجموعهٔ نامتناهی از ارزش‌های صدق، اغلب
بررسی زیرمجموعه‌های زیر نیز سودمند است:}
V_m & = \Setabs{\frac{n}{m-1}}{n \in \Nat \text{ و } n<m}
\intertext{برای نمونه، $V_5$ مجموعه‌ای با پنج ارزش صدق هم‌فاصله است:}
V_5 & = \{0, \frac{1}{4}, \frac{1}{2}, \frac{3}{4}, 1\}.
\end{align*}
در منطق‌های مبتنی بر این مجموعه‌های ارزش صدق، معمولاً فقط~$1$ ممتاز
است؛ یعنی~$V^+ = \{1\}$. به بیان دیگر، $1$ را در نقش صدق مطلق و~$0$
را در نقش کذب مطلق می‌گذاریم، اما~!!{formula}ها می‌توانند هر مقدار
میانی در~$V$ را بگیرند.

همچنین می‌توان مجموعهٔ~$V_{[0,1]} = [0,1]$ از همهٔ اعداد
\emph{حقیقی} میان~$0$ و~$1$، یا زیرمجموعه‌های نامتناهی دیگری از
$[0,1]$ را در نظر گرفت. منطق‌های دارای این مجموعهٔ ارزش صدق اغلب
\emph{فازی} نامیده می‌شوند.


\end{document}
